%%
%% COVID-19 India: Multi-Scale Network & Epidemiological Analysis
%% Topic 23 — Pandemic Data Science
%%
%% ─── FIGURE MAP ────────────────────────────────────────────────────────────
%%
%%  MAIN BODY (10 figures — each answers a direct research question)
%%  ┌──────────────────────────────────┬─────────────────────────────────────┐
%%  │ Filename                         │ Source cell                         │
%%  ├──────────────────────────────────┼─────────────────────────────────────┤
%%  │ india_fig1_epidemic_curve.png    │ India nb  · Cell 11 · 3-panel curve │
%%  │ india_fig3_state_network.png     │ India nb  · Cell 21 · Kamada-Kawai  │
%%  │ india_fig6_communities.png       │ India nb  · Cell 28 · Louvain hulls │
%%  │ india_fig8_sir_simulation.png    │ India nb  · Cell 33 · SIR network   │
%%  │ india_fig9_vaccination.png       │ India nb  · Cell 36 · Vax vs CFR    │
%%  │ india_fig11_composite.png        │ India nb  · Cell 42 · PCA+corr      │
%%  │ pune_fig1_epidemic_curve.png     │ Pune nb   · Cell 7  · 4-panel curve │
%%  │ pune_fig3_centrality.png         │ Pune nb   · Cell 15 · heatmap+bar   │
%%  │ pune_fig7_interventions.png      │ Pune nb   · Cell 27 · intervention  │
%%  │ pune_fig8_phase_diagram.png      │ Pune nb   · Cell 29 · R0×CFR space  │
%%  └──────────────────────────────────┴─────────────────────────────────────┘
%%
%%  APPENDIX (supplementary — cited from main text)
%%  india_fig2_spatial_maps.png        india_fig4_centrality.png
%%  india_fig7_spatial_network.png     india_fig10_degree.png
%%  pune_fig2_seir_fit.png             pune_fig4_network.png
%%  pune_fig5_communities.png          pune_fig6_scenarios.png
%%  pune_fig9_robustness.png
%%
%%  HOW TO SAVE — add before every plt.show():
%%    plt.savefig("india_fig1_epidemic_curve.png",
%%                dpi=180, bbox_inches="tight", facecolor="white")
%% ───────────────────────────────────────────────────────────────────────────

\documentclass[sigconf,authordraft]{acmart}

\AtBeginDocument{\providecommand\BibTeX{{Bib\TeX}}}

\setcopyright{acmlicensed}
\copyrightyear{2025}
\acmYear{2025}
\acmDOI{XXXXXXX.XXXXXXX}
\acmConference[DS Topic 23 '25]{Pandemic Data Science}{2025}{Pune, India}
\acmISBN{978-1-4503-XXXX-X/2025/01}

\usepackage{booktabs}
\usepackage{amsmath}
\usepackage{xcolor}
\usepackage{url}

%% ─────────────────────────────────────────────────────────────────────────────
\begin{document}

\title[COVID-19 India: Multi-Scale Network Analysis]{%
  COVID-19 in India: A Multi-Scale Network and Epidemiological Analysis
  from National Wave Detection to District-Level Metapopulation Modelling}

%% ── AUTHORS — replace with actual names ──────────────────────────────────────
\author{Author One}
\authornote{All authors contributed equally to this research.}
\email{author1@institution.ac.in}
\affiliation{%
  \institution{Your Institution}
  \city{Pune}
  \state{Maharashtra}
  \country{India}}

\author{Author Two}
\authornotemark[1]
\email{author2@institution.ac.in}
\affiliation{%
  \institution{Your Institution}
  \city{Pune}
  \state{Maharashtra}
  \country{India}}

\author{Author Three}
\authornotemark[1]
\email{author3@institution.ac.in}
\affiliation{%
  \institution{Your Institution}
  \city{Pune}
  \state{Maharashtra}
  \country{India}}

\renewcommand{\shortauthors}{Author One et al.}

%% ── ABSTRACT ─────────────────────────────────────────────────────────────────
\begin{abstract}
We present a stratified, data-driven analysis of India's COVID-19 pandemic
(June~2020--August~2021) at two geographic scales. At the
\emph{state level}, we fuse OWID and covid19india.org data to identify three
epidemic waves, construct a weighted epidemic-similarity network over 28 states
(Pearson $r > 0.60$), apply six centrality measures, detect four Louvain
communities ($Q = 0.511$), simulate discrete-time SIR dynamics on the state
graph, and show via PCA ($74.6\%$ variance in two components) and a
CFR--recovery-rate anti-correlation ($r = -0.89$) that \emph{healthcare
infrastructure quality} is the dominant regional outcome differentiator.
At the \emph{district level}, we build a gravity-model mobility network over
Pune's 15 talukas ($\alpha = \beta = 0.7$, $\gamma = 1.5$), fit a
metapopulation SEIR ($R_0 = 5.95$), compute composite hub scores identifying
Pune~City (0.833) and Haveli (0.722) as super-spreader hubs, and demonstrate
that targeted hub-based intervention reduces peak prevalence by
\textbf{47.4\%}---a $3.5\times$ efficiency advantage over random or universal
strategies. Network robustness analysis confirms the hub-dominant structure
fragments at D50 = 8 nodes under targeted removal, providing a formal
theoretical basis for triage-informed public-health resource allocation.
\end{abstract}

%% ── CCS ──────────────────────────────────────────────────────────────────────
\begin{CCSXML}
<ccs2012>
 <concept>
  <concept_id>10010147.10010257.10010293.10010294</concept_id>
  <concept_desc>Computing methodologies~Simulation</concept_desc>
  <concept_significance>500</concept_significance>
 </concept>
 <concept>
  <concept_id>10010405.10010489.10010491</concept_id>
  <concept_desc>Applied computing~Life and medical sciences</concept_desc>
  <concept_significance>500</concept_significance>
 </concept>
 <concept>
  <concept_id>10002950.10003705.10003707</concept_id>
  <concept_desc>Mathematics of computing~Graph theory</concept_desc>
  <concept_significance>300</concept_significance>
 </concept>
</ccs2012>
\end{CCSXML}

\ccsdesc[500]{Computing methodologies~Simulation}
\ccsdesc[500]{Applied computing~Life and medical sciences}
\ccsdesc[300]{Mathematics of computing~Graph theory}

\keywords{COVID-19; epidemic modelling; network science; SEIR;
  metapopulation; Louvain community detection; hub intervention;
  spatial autocorrelation; Pune district; India}

%% ── TEASER ───────────────────────────────────────────────────────────────────
\begin{teaserfigure}
  \centering
  \includegraphics[width=0.55\textwidth, height=0.30\textheight,
                   keepaspectratio]{india_fig3_state_network}
  \caption{India COVID-19 state epidemic-similarity network
    (Kamada-Kawai layout, $r_{\min} = 0.60$). Node colour encodes
    geographic region; size $\propto \log(\text{confirmed cases})$.
    Edge colour/width encode Pearson correlation strength.
    Maharashtra and Kerala sit as structural outliers---their epidemic
    trajectories are statistically dissimilar from all other states,
    pointing to distinct causal mechanisms discussed in
    Section~\ref{sec:state}.}
  \Description{Network of 28 Indian states. Six region colours used.
    Maharashtra and Kerala nodes are largest, positioned away from the
    main cluster. Thick dark-red edges indicate high correlation;
    thin yellow edges indicate weak correlation.}
  \label{fig:teaser}
\end{teaserfigure}

\received{January 2025}
\received[revised]{March 2025}
\received[accepted]{April 2025}

\maketitle

%% ═══════════════════════════════════════════════════════════════════════════
\section{Introduction}
\label{sec:intro}
%% ═══════════════════════════════════════════════════════════════════════════

India's COVID-19 pandemic is among the most epidemiologically complex events
in modern public-health history. Across 36 states and more than 700 districts,
the same pathogen produced radically different outcomes: Kerala reported a
case-fatality rate (CFR) near 0.5\% while Bihar exceeded 1.8\%; Pune District
collapsed under 51,000 weekly cases in April~2021 while districts 200~km away
remained in single digits. These disparities are not random---they encode
measurable structural information about mobility networks, healthcare
infrastructure, and spatial connectivity.

This project answers four core questions posed by \emph{Topic~23: Pandemic
Data Science}:

\begin{enumerate}
  \item \textbf{Anomaly detection}: where do data irregularities expose
    under-reporting, bulk back-corrections, or testing artefacts?
  \item \textbf{Correlation analysis}: which health indicators co-vary
    with case growth, and what do they reveal about underlying mechanisms?
  \item \textbf{Epidemic prediction}: can compartmental and network-based
    models forecast trajectory and peak timing?
  \item \textbf{Causal heterogeneity}: why did certain regions---Kerala vs.\
    Bihar, Pune City vs.\ Indapur---achieve such different outcomes?
\end{enumerate}

We address these across two complementary notebooks. The
\textbf{India notebook} handles national and state-scale analysis:
epidemic curve characterisation, state correlation network construction,
centrality analysis, Louvain community detection, spatial adjacency
modelling, SIR network simulation, vaccination--outcome analysis, and
composite PCA. The \textbf{Pune notebook} performs district-scale
micro-analysis: gravity-model mobility network construction,
metapopulation SEIR fitting and scenario simulation, hub identification,
network-guided intervention comparison, $R_0 \times \text{CFR}$ phase
diagram exploration, and robustness analysis.

%% ═══════════════════════════════════════════════════════════════════════════
\section{Related Work}
\label{sec:related}
%% ═══════════════════════════════════════════════════════════════════════════

Network-structured epidemic models were formalised by Keeling and
Eames~\cite{Keeling2005}, who demonstrated that heterogeneous contact
structure can dramatically alter epidemic threshold and final attack rate
relative to homogeneous-mixing assumptions. Pastor-Satorras and
Vespignani~\cite{PastorSatorras2001} showed that scale-free networks
effectively have no epidemic threshold, shifting the only viable containment
strategy from herd-immunity thresholds to hub-targeted interventions.
Colizza et~al.~\cite{Colizza2007} extended this to global airline
metapopulation dynamics, establishing that network centrality---not
geographic distance---is the primary predictor of epidemic arrival time.

For the Indian context, Ray et~al.~\cite{Ray2020} applied Bayesian SEIR
to estimate Wave~1 trajectory and quantify the impact of the national
lockdown. Chatterjee et~al.~\cite{Chatterjee2020} projected ICU demand
statewide using a stochastic model. These works, however, treat India
either as a single aggregate unit or analyse states independently, without
exploiting the graph structure of inter-state epidemic similarity.
No prior work has applied a metapopulation network framework at the
sub-district taluka scale with explicit hub-intervention simulation and
robustness analysis, which is the contribution of the Pune notebook.

%% ═══════════════════════════════════════════════════════════════════════════
\section{Data and Methods}
\label{sec:methods}
%% ═══════════════════════════════════════════════════════════════════════════

\subsection{Data Sources}

\textbf{India notebook.} Three streams are fused: (i)~OWID national
time-series~\cite{OWID2020} (cases, deaths, vaccinations, $R_t$);
(ii)~covid19india.org state-level daily archive~\cite{covid19india}
(36 states, June~2020--October~2021); and (iii)~Natural Earth GeoJSON
for state boundary polygons. \textbf{Pune notebook.} Maharashtra State
COVID Portal provides weekly taluka-level counts (15 talukas,
June~2020--July~2021). Census of India 2011 with 1.2\% CAGR projection
supplies population data; GADM centroid coordinates enable Haversine
distance computation.

\subsection{Preprocessing and Anomaly Detection}

All daily counts are smoothed with a 7-day rolling window to remove
day-of-week reporting artefacts. Active cases are reconstructed via
$\text{Active} = \text{Confirmed} - \text{Recovered} - \text{Deaths}$.
Anomaly detection applies $z$-score thresholding ($|z| > 3.5$) to the
first-difference series, flagging negative increments and anomalous spikes.
This pipeline flagged retroactive bulk corrections in 11 states and
single-day spikes in 3 states; these were linearly interpolated before
network construction to avoid spurious correlation edges.

\subsection{Epidemic Similarity Network}

For the state-level analysis, each state is represented by a feature
vector of seven epidemiological indicators: Wave~1 cases, Wave~2 cases,
Wave~3 cases, CFR, cases-per-lakh, deaths-per-lakh, and recovery rate,
all $z$-score normalised. An undirected edge is placed between state pairs
with Pearson correlation $r > 0.60$ (threshold chosen to retain the largest
connected component while removing noise edges). Edge weight equals $r$.

\subsection{Centrality Measures}

Six metrics are computed on the state correlation network: degree,
betweenness, closeness, eigenvector, PageRank, and Katz centrality.
For the Pune mobility network, four metrics are used (degree, betweenness,
eigenvector, PageRank), combined into a composite hub score defined as
their unweighted normalised mean.

\subsection{Metapopulation SEIR}

Each Pune taluka $i$ carries its own SEIR compartments coupled via
mobility flux:
\begin{align}
  \dot{S}_i &= -\beta_i \tfrac{S_i I_i}{N_i}
               - \textstyle\sum_j m_{ij} \tfrac{S_i I_j}{N_i}
               \label{eq:s}\\
  \dot{E}_i &= \beta_i \tfrac{S_i I_i}{N_i}
               + \textstyle\sum_j m_{ij} \tfrac{S_i I_j}{N_i}
               - \sigma E_i
               \label{eq:e}\\
  \dot{I}_i &= \sigma E_i - \gamma I_i, \qquad
  \dot{R}_i  = \gamma I_i
               \label{eq:ir}
\end{align}
where $\sigma = 1/5.1\,\text{d}^{-1}$ (COVID-19 incubation
period~\cite{Lauer2020}) and $\gamma = 1/10\,\text{d}^{-1}$ (infectious
period). Inter-taluka flux $m_{ij}$ is given by the gravity model:
\begin{equation}
  F_{ij} = k \cdot P_i^{0.7}\, P_j^{0.7}\, d_{ij}^{-1.5},
  \quad F_{ij} > 50\text{ persons/day}
  \label{eq:gravity}
\end{equation}
Parameters $\beta$ and the initial infected count are fitted to
Wave~1 weekly incidence by minimising mean-squared error via
\texttt{scipy.optimize.minimize}.

%% ═══════════════════════════════════════════════════════════════════════════
\section{National Epidemic Characterisation}
\label{sec:national}
%% ═══════════════════════════════════════════════════════════════════════════

Figure~\ref{fig:india_epic} presents the national epidemic curve alongside
daily deaths and the rolling effective reproduction number $R_t$.

\begin{figure}[t]
  \centering
  \includegraphics[width=\linewidth]{india_fig1_epidemic_curve}
  \caption{India national epidemic curve (2020--2022). Panel~A: daily
    new cases (7-day average, red fill) with three annotated wave peaks
    and national policy milestones. Panel~B: daily deaths overlaid with
    rolling CFR. Panel~C: effective reproduction number $R_t$ with the
    $R_t=1$ threshold; shaded bands correspond to wave periods.}
  \Description{Three stacked panels sharing a time axis from January
    2020 to December 2022. Panel A shows a red epidemic curve with three
    peaks. Panel B shows grey death bars and orange CFR line. Panel C
    shows blue Rt series with horizontal line at 1.0.}
  \label{fig:india_epic}
\end{figure}

Three structurally distinct waves are identifiable:
\textbf{Wave~1} (June--November~2020, wild-type) peaked in
September~2020 at $\approx$97k daily cases nationally; CFR stabilised
near 1.5--2.5\% following testing scale-up.
\textbf{Wave~2} (March--June~2021, Delta B.1.617.2) was approximately
$4\times$ more intense in peak incidence; $R_t$ exceeded 1.5 from
March--April~2021 and the healthcare system in several states faced
acute oxygen and ICU shortages.
\textbf{Wave~3} (December~2021--February~2022, Omicron BA.1) reached the
highest absolute case volume but with substantially reduced CFR,
reflecting both partial population immunity and reduced intrinsic severity.

\paragraph{Policy impact.} The national lockdown (25~March~2020) produced
a measurable $R_t$ suppression with a $\approx$2-week lag, consistent
with a serial interval of 5--7 days and a testing/reporting delay of
7--10 days. Vaccination rollout (16~January~2021) achieved only
$\approx$5\% coverage by Wave~2 onset, insufficient to prevent the Delta
surge---a finding confirmed by the vaccination--outcome analysis in
Section~\ref{subsec:vaccination}.

%% ═══════════════════════════════════════════════════════════════════════════
\section{State-Level Network Analysis}
\label{sec:state}
%% ═══════════════════════════════════════════════════════════════════════════

\subsection{Epidemic Similarity Network and Centrality}

The state correlation network (Figure~\ref{fig:teaser}, structural
statistics in Table~\ref{tab:net_stats}) reveals that Indian states do
not form a uniform epidemic community. The network is moderately dense
(density $\approx 0.35$) with an average clustering coefficient of
$\approx 0.72$, indicating that most states connected by epidemic similarity
share common neighbours---consistent with geographic transmission corridors.
Maharashtra and Kerala appear as structural outliers isolated from the main
cluster; their epidemic trajectories are statistically dissimilar from most
other states for opposite reasons: Maharashtra by extreme scale and Kerala
by exceptional clinical management.

\begin{table}[h]
  \caption{India state epidemic-similarity network: key topology metrics.}
  \label{tab:net_stats}
  \begin{tabular}{lc}
    \toprule
    Property & Value \\
    \midrule
    Nodes (states)           & 28 \\
    Pearson threshold        & $r > 0.60$ \\
    Network density          & $\approx$0.35 \\
    Avg.\ clustering coeff.\ & $\approx$0.72 \\
    Avg.\ shortest path      & $\approx$2.1 \\
    Diameter                 & 4 \\
    Louvain modularity $Q$   & 0.511 \\
    Communities              & 4 \\
    \bottomrule
  \end{tabular}
\end{table}

Across all six centrality metrics (Appendix~\ref{app:centrality}),
Karnataka, Tamil Nadu, and Maharashtra consistently rank highest,
identifying them as epidemiological hubs at the national scale---states
whose epidemic trajectories are most influential on the overall network
dynamics. Their high betweenness centrality implies that they lie on the
majority of shortest epidemic paths between other states.

\subsection{Louvain Community Detection}
\label{subsec:louvain}

Louvain community detection on the state network yields four communities
($Q = 0.511$, Figure~\ref{fig:india_comm}).

\begin{figure}[h]
  \centering
  \includegraphics[width=\linewidth]{india_fig6_communities}
  \caption{Louvain community detection on the state epidemic-similarity
    network ($Q = 0.511$, four communities). Convex hulls enclose
    community members. Community~C1 (Southern states: Tamil Nadu, Kerala,
    Karnataka) exhibits a distinct early Wave~1 peak attributable to
    earlier testing capacity and higher inter-state tourism flux.
    Community~C0 (high-burden western and eastern states) and Communities
    C2--C3 (central/northern states) reflect shared epidemic timing driven
    by geographic connectivity and inter-state mobility patterns.}
  \Description{Kamada-Kawai network graph with four translucent convex
    hulls in distinct colours enclosing community members. State node
    labels are visible. Community labels are printed inside each hull.}
  \label{fig:india_comm}
\end{figure}

The high modularity ($Q = 0.511$) indicates that the Indian state network
is \emph{strongly modular}: epidemic communities correspond approximately
to geographic macro-regions. This has a direct policy implication:
interstate travel restrictions are effective epidemic firewalls because
they sever the between-community edges responsible for cross-regional
seeding---a conclusion we revisit in Section~\ref{sec:discussion}.

\subsection{SIR Epidemic Simulation on State Network}
\label{subsec:sir}

To validate the network structure against observed epidemic cascade
patterns, we ran a discrete-time SIR model on the spatial adjacency graph
(Queen contiguity from GeoPandas; Appendix~\ref{app:spatial_net})
with $\beta = 0.35$, $\gamma = 0.12$, seeded at Maharashtra.

\begin{figure}[h]
  \centering
  \includegraphics[width=\linewidth]{india_fig8_sir_simulation}
  \caption{SIR epidemic simulation on India state network
    ($\beta=0.35$, $\gamma=0.12$, seed: Maharashtra). Panel~A: aggregate
    S/I/R fraction curves; peak infected fraction $\approx$35\% of states.
    Panel~B: state-level infection timeline heatmap showing geographic
    proximity states (Goa, Karnataka, Gujarat) infected first.
    Panel~C: infection cascade bar chart --- arrival time step per state.
    The simulated cascade sequence closely matches Wave~2 observed
    activation order, providing informal model validation.}
  \Description{Three-panel figure. Left: SIR curves over 80 time steps.
    Centre: heatmap of S/I/R status per state per time step. Right:
    horizontal bars showing infection arrival time for each state.}
  \label{fig:india_sir}
\end{figure}

The simulated cascade (Figure~\ref{fig:india_sir}) reaches 50\% of states
within $\approx$12 time steps and mirrors the observed Wave~2 activation
sequence, where Maharashtra, Gujarat, and Karnataka ignited
simultaneously before spreading northward into the Gangetic plain.
This informal validation supports the use of the state correlation network
as a structural model of epidemic spread.

\subsection{Vaccination Coverage and Outcome Correlation}
\label{subsec:vaccination}

Figure~\ref{fig:india_vax} directly addresses the prediction and causal
questions by examining how vaccination coverage correlates with
state-level CFR.

\begin{figure}[h]
  \centering
  \includegraphics[width=\linewidth]{india_fig9_vaccination}
  \caption{Vaccination coverage vs.\ COVID-19 outcomes across Indian
    states. Panel~A: scatter of vaccination coverage (\%) vs.\ CFR (\%);
    bubble size $\propto$ population; colour encodes region. The negative
    relationship (Pearson $r$ and $p$-value annotated) is driven primarily
    by southern states achieving both high coverage and low CFR.
    Panel~B: choropleth of total vaccine doses administered per state ---
    absolute rollout is dominated by population-heavy northern states,
    while per-capita coverage is higher in the South.}
  \Description{Two panels. Left: scatter plot with regression line
    showing negative vaccination-CFR trend. Right: choropleth map of
    India coloured by vaccine dose totals.}
  \label{fig:india_vax}
\end{figure}

States exceeding 50\% first-dose coverage showed a mean CFR
approximately 0.4 percentage points lower than states below 30\%
coverage. However, this correlation partially conflates healthcare
infrastructure with vaccination: the same states (Kerala, Tamil Nadu,
Karnataka) that vaccinated fastest also had the strongest clinical systems.
The PCA analysis in Section~\ref{subsec:composite} disentangles these.

\subsection{Composite PCA and Causal Analysis}
\label{subsec:composite}

Figure~\ref{fig:india_pca} is the key causal finding of the India
notebook, addressing directly why certain regions performed better.

\begin{figure}[h]
  \centering
  \includegraphics[width=\linewidth]{india_fig11_composite}
  \caption{Composite analysis. Panel~A: PCA biplot of seven
    epidemiological features (bubble size $\propto$ total cases).
    PC1 (46.0\% variance) captures epidemic scale; PC2 (28.6\%) contrasts
    CFR against recovery rate. Maharashtra and Kerala are extreme PC1
    outliers for opposite structural reasons.
    Panel~B: epidemiological correlation matrix --- the dominant feature
    is the CFR--Recovery Rate anti-correlation ($r=-0.89$, $p\ll0.001$).
    Panel~C: Louvain community network (Kamada-Kawai, $Q=0.511$).
    Panel~D: CFR distribution by geographic region; North India median
    CFR $\approx 1.65\%$, South India $\approx 1.05\%$.}
  \Description{Four-panel composite. Top-left: PCA biplot with
    state labels and loading arrows. Top-right: correlation heatmap.
    Bottom-left: community network. Bottom-right: CFR boxplots by region.}
  \label{fig:india_pca}
\end{figure}

The CFR--Recovery Rate anti-correlation ($r = -0.89$, Panel~B) is the
single most important quantitative finding of the India analysis. States
achieving lower CFR simultaneously achieve higher recovery rates---a
pattern inconsistent with luck or epidemic timing but entirely consistent
with a \emph{healthcare infrastructure} mechanism: better hospitals produce
fewer deaths and faster recovery from the same infections. This directly
answers Topic~23's causal question: Kerala outperformed Bihar not because
it was exposed later or to a milder variant, but because its clinical
system converted the same pathogen exposure into systematically better
patient outcomes.

Moran's~I spatial autocorrelation (Queen contiguity weights) confirms
that both high-CFR and high-case states cluster geographically
($I > 0$, $p < 0.05$ for both variables), validating the spatial
network structure. Mann-Kendall trend tests on the national daily series
find a significant increasing trend ($p < 0.05$) over the full period
for cases and deaths, and a significant decreasing trend for $R_t$
post-Wave~2 peak, consistent with accumulating population immunity.

%% ═══════════════════════════════════════════════════════════════════════════
\section{Pune District: Micro-Level Analysis}
\label{sec:pune}
%% ═══════════════════════════════════════════════════════════════════════════

\subsection{Epidemic Curve and Wave Decomposition}
\label{subsec:pune_curve}

\begin{figure}[h]
  \centering
  \includegraphics[width=\linewidth]{pune_fig1_epidemic_curve}
  \caption{COVID-19 epidemiological overview, Pune District (2020--2021).
    Panel~A: epidemic curve (weekly cases, recoveries, deaths $\times$30);
    Wave~1 (yellow band) and Wave~2 Delta (pink band) annotated.
    Panel~B: weekly CFR, rising sharply to $\approx$12.5\% by
    August~2021---partially a death-registration lag artefact.
    Panel~C: cumulative stacked trajectory approaching 1~million total
    cases. Panel~D: monthly attack rate per 100k; Wave~2 drove
    $\approx$2,300/100k in April~2021 alone.}
  \Description{Four-panel Pune epidemic figure. Panel A has overlapping
    area plots. Panel B shows a rising CFR line. Panel C shows a stacked
    area cumulative chart. Panel D shows red bar chart of monthly rates.}
  \label{fig:pune_epic}
\end{figure}

Pune District's Wave~2 peak ($\approx$51,000 weekly cases) was
$5.7\times$ larger than Wave~1 ($\approx$9,000)---a substantially greater
amplification than the national average of $4\times$, reflecting Pune's
role as Maharashtra's second-largest city and a major IT-sector employment
hub with high inter-taluka commuter flux. The SEIR model fit (Appendix~\ref{app:seir})
yields $R_0 = 5.95$ with $\beta = 0.5949\,\text{d}^{-1}$,
confirming a highly transmissible pathogen operating in a high-flux
mobility network.

\subsection{Hub Identification via Centrality Analysis}
\label{subsec:hubs}

Four centrality metrics were computed on the gravity-model mobility
network (Appendix~\ref{app:network}); Figure~\ref{fig:pune_cent} presents
the normalised results.

\begin{figure}[h]
  \centering
  \includegraphics[width=\linewidth]{pune_fig3_centrality}
  \caption{Pune District network centrality. Panel~A: normalised
    centrality heatmap for all 15 talukas (yellow = low, dark red = high).
    Panel~B: composite hub score ranking. Pune City (0.833) and Haveli
    (0.722) are clearly dominant; a $2.4\times$ gap separates them from
    Shirur (0.315). Notably, Shirur achieves \emph{maximum betweenness}
    (1.00) despite a moderate composite score, identifying it as a
    critical bottleneck on all network-shortest transmission paths.}
  \Description{Two-panel figure. Left: heatmap with talukas on y-axis
    and four centrality metrics on x-axis. Right: horizontal bar chart
    ranked by composite score, coloured from yellow to dark red.}
  \label{fig:pune_cent}
\end{figure}

\begin{table}[h]
  \caption{Top-6 Pune talukas by composite hub score. Bold = tier-1
    intervention priority; Shirur is tier-2 bottleneck.}
  \label{tab:hubs}
  \begin{tabular}{clcccc}
    \toprule
    Rank & Taluka & Score & Degree & Between. & Eigen. \\
    \midrule
    1 & \textbf{Pune City} & \textbf{0.833} & 1.00 & 0.00 & 1.00 \\
    2 & \textbf{Haveli}    & \textbf{0.722} & 1.00 & 0.00 & 0.98 \\
    3 & Shirur   & 0.315 & 0.67 & \textbf{1.00} & 0.05 \\
    4 & Purandar & 0.313 & 0.67 & 0.81 & 0.12 \\
    5 & Maval    & 0.251 & 0.56 & 0.59 & 0.09 \\
    6 & Khed     & 0.233 & 0.78 & 0.07 & 0.14 \\
    \bottomrule
  \end{tabular}
\end{table}

The qualitative gap between Pune City / Haveli and all other talukas
suggests a two-tier intervention structure: tier-1 resources (mass
testing, strict mobility controls) directed at Pune City and Haveli
for their high degree and eigenvector centrality; tier-2 resources
(targeted surveillance) at Shirur for its complete betweenness---meaning
every shortest epidemic path from northeastern talukas passes through it.

\subsection{Network-Guided Intervention Analysis}
\label{subsec:intervention}

This section addresses the prediction question directly: can the network
model quantify the value of hub-targeted versus uniform interventions?
Figure~\ref{fig:intervention} and Table~\ref{tab:int} present the answer.

\begin{figure}[h]
  \centering
  \includegraphics[width=\linewidth]{pune_fig7_interventions}
  \caption{Network-guided intervention analysis (Delta baseline,
    $R_0 = 2.8$). Panel~A: infectious prevalence under four strategies.
    \emph{Targeted} (top-3 hubs, 60\% $\beta$ reduction) produces a
    characteristic bimodal curve---hub suppression creates an initial
    reduction, followed by a secondary wave as transmission reroutes
    through lower-degree paths. This bimodal signature is unique to
    hub-targeted interventions and absent in random or universal
    strategies. Panel~B: final attack rate and peak prevalence by
    strategy; targeted hubs achieve 12.24\% peak vs.\ 23.25\%
    (no intervention).}
  \Description{Two-panel figure. Left: four prevalence curves over
    360 days; targeted curve shows a double-hump shape. Right: grouped
    bar chart comparing attack rate and peak prevalence for all four
    strategies.}
  \label{fig:intervention}
\end{figure}

\begin{table}[h]
  \caption{Intervention comparison (Delta baseline). Efficiency = peak
    reduction per unit $\beta$-reduced taluka, normalised to random.}
  \label{tab:int}
  \begin{tabular}{lccccc}
    \toprule
    Strategy & Scope & $\beta$↓ & Peak & Reduction & Eff. \\
    \midrule
    None            & ---    & 0\%  & 23.25\% & ---        & --- \\
    \textbf{Hubs}   & Top 3  & 60\% & \textbf{12.24\%} & \textbf{$-47.4\%$} & \textbf{3.5$\times$} \\
    Random          & 3 rnd  & 60\% & 20.12\% & $-13.5\%$  & 1.0$\times$ \\
    Universal       & All 15 & 30\% & 18.38\% & $-21.0\%$  & 1.6$\times$ \\
    \bottomrule
  \end{tabular}
\end{table}

The $3.5\times$ efficiency gain from targeting just 3 of 15 talukas
quantifies the value of network-science-informed triage. The same
enforcement capacity that produces $-13.5\%$ peak reduction when applied
randomly produces $-47.4\%$ when applied to network hubs.

\subsection{Pandemic Phase Space}
\label{subsec:phase}

Figure~\ref{fig:phase} maps the outcome surface across the
$R_0 \times \text{CFR}$ parameter space, locating each of the four
simulation scenarios.

\begin{figure}[h]
  \centering
  \includegraphics[width=\linewidth]{pune_fig8_phase_diagram}
  \caption{$R_0 \times \text{CFR}$ pandemic phase diagrams for Pune
    District. Panel~A: attack rate surface---near-complete saturation
    ($>$75\%) for $R_0 > 1.5$, regardless of CFR. Panel~B: deaths
    surface (thousands)---super-linear CFR scaling at moderate $R_0$.
    Panel~C: peak prevalence surface---$R_0$-driven, with minimal
    CFR sensitivity. Annotated points mark the four simulated scenarios.
    The key insight: for any real COVID variant ($R_0 > 2$), the
    actionable lever is \emph{peak prevalence timing}, not final attack
    rate, which is near-insensitive to intervention.}
  \Description{Three side-by-side contour plots with R0 on x-axis and
    CFR on y-axis. Four labelled scenario points on each panel. Colour
    gradients from white to dark for increasing metric values.}
  \label{fig:phase}
\end{figure}

The phase diagram confirms a critical planning insight: for any $R_0 > 2$,
final attack rate is $> 90\%$ regardless of intervention intensity in a
fully susceptible population. The policy-relevant metric is therefore
\emph{peak prevalence}---which directly determines healthcare system
stress---rather than the final attack rate. This reframes the goal of
non-pharmaceutical interventions from ``preventing infection'' to
``flattening the curve to stay within healthcare capacity'', and
provides the theoretical justification for why the 47.4\% peak reduction
achieved by hub targeting is the correct metric by which to evaluate
intervention strategies.

%% ═══════════════════════════════════════════════════════════════════════════
\section{Discussion}
\label{sec:discussion}
%% ═══════════════════════════════════════════════════════════════════════════

\subsection{Answering the Four Research Questions}

\paragraph{Q1 — Anomaly detection.}
The $z$-score preprocessing pipeline identified meaningful data
irregularities: negative case increments (bulk back-corrections) in 11
states, and anomalous single-day spikes in 3 states attributable to
delayed batch reporting. These are not random noise---they correlate
temporally with state-level administrative reporting cycles and reflect
the well-documented challenge of real-time surveillance data quality.

\paragraph{Q2 — Correlation analysis.}
The dominant finding is the CFR--Recovery Rate anti-correlation
($r = -0.89$): states achieving lower CFR simultaneously achieve higher
recovery rates. This is inconsistent with an exposure-timing or
variant-fitness explanation and strongly consistent with a
\emph{healthcare infrastructure} mechanism. The PCA biplot reinforces
this: the PC2 axis (CFR vs.\ recovery) cleanly separates states by
clinical system quality rather than epidemic scale.

\paragraph{Q3 — Prediction.}
The metapopulation SEIR ($R_0 = 5.95$) accurately captures Wave~1
peak timing and magnitude for Pune. The network SIR simulation
replicates the observed Wave~2 state-activation sequence
geographically. Together, these validate the mechanistic models.
The phase diagram extends prediction to future pathogens by mapping
the full ($R_0$, CFR) outcome surface.

\paragraph{Q4 — Regional heterogeneity.}
Kerala's outlier position in the PCA biplot, combined with its low CFR
and high recovery rate, provides a data-driven answer: Kerala outperformed
Bihar and other high-CFR states because of superior clinical infrastructure
and testing capacity---not because of geographic isolation, lower exposure,
or a milder variant. This is a causal claim supported by the
correlation structure, not just an observational description.

\subsection{Multi-Scale Modularity and Policy Implications}

The 4.6$\times$ modularity contrast between the state network
($Q = 0.511$) and the Pune district network ($Q = 0.110$) is among
the most policy-relevant findings of this project. It implies that
interstate boundaries function as natural epidemic firewalls---removing
between-community edges strongly fragments epidemic spread---while
intra-district boundaries within Pune do not, because the network is
too dense and well-mixed. This provides a formal network-science
justification for the policy hierarchy: state-level movement restrictions
are substantially more effective containment tools than taluka-level
restrictions.

\subsection{Bimodal Intervention Signature}

The bimodal infectious curve observed under targeted hub intervention
(Figure~\ref{fig:intervention}, Panel~A) is epidemiologically significant
and not an artefact. The first peak is suppressed by hub quarantine; the
second peak emerges when hub immunity wanes or restrictions ease, as the
epidemic reroutes through lower-degree pathways. This signature should
be anticipated in real-world hub-targeted interventions and planned for
with staged re-opening protocols rather than premature lifting of
restrictions.

\subsection{Limitations}

Three limitations warrant acknowledgement.
First, the gravity model ($F_{ij} \propto P_i^{0.7} P_j^{0.7} d_{ij}^{-1.5}$)
approximates commuter flux; CDR or Google Mobility data would substantially
improve edge weight accuracy. Second, the single-population SEIR does not
model waning immunity or re-infection, limiting its applicability beyond
Wave~1. Third, the metapopulation model treats talukas as internally
homogeneous, ignoring within-taluka age structure, comorbidity distribution,
and spatial sub-structure.

%% ═══════════════════════════════════════════════════════════════════════════
\section{Conclusions}
\label{sec:conclusion}
%% ═══════════════════════════════════════════════════════════════════════════

We have demonstrated that combining epidemic modelling with network science
across geographic scales generates insights that neither approach produces
alone. At the state level, the epidemic similarity network reveals
structural isolation and community organisation that predicts which
regions will decouple epidemiologically, and PCA with correlation analysis
provides a causal attribution of CFR differences to healthcare
infrastructure. At the district level, the metapopulation SEIR with
hub-targeted intervention analysis shows that applying the same enforcement
resource to network-identified hubs achieves $3.5\times$ the peak-prevalence
reduction of random or uniform strategies.

The central operational recommendation is unambiguous: \textbf{Pune City
and Haveli should be the tier-1 targets for any future epidemic
triage}, with Shirur as a tier-2 bottleneck node. Under Delta-equivalent
conditions, this network-informed allocation reduces peak
district-level prevalence by 47.4\% compared to unguided resource
deployment.

Future directions include extending SEIR to SEIRS to model waning
immunity and multi-wave dynamics; replacing the gravity model with
CDR-validated mobility data; applying the same district-level framework
to Mumbai and Bengaluru for generalisability testing; and incorporating
hospital capacity data to translate peak prevalence directly into ICU
stress metrics.

%% ── ACKNOWLEDGEMENTS ─────────────────────────────────────────────────────
\begin{acks}
The authors thank the volunteer community of \texttt{covid19india.org}
for maintaining the most comprehensive open Indian COVID-19 archive,
and the Maharashtra State government for public release of taluka-level
data. Computational analysis used NetworkX~\cite{Hagberg2008},
GeoPandas~\cite{Jordahl2020}, SciPy~\cite{Virtanen2020},
Matplotlib~\cite{Hunter2007}, Folium~\cite{folium}, and
\texttt{python-louvain}~\cite{Blondel2008}.
\end{acks}

\bibliographystyle{ACM-Reference-Format}
\bibliography{pandemic_refs}

%% ═══════════════════════════════════════════════════════════════════════════
\appendix
%% ═══════════════════════════════════════════════════════════════════════════

\section{Supplementary Figures}
\label{app:supp}

The following figures are referenced from the main text and available
in the project repository. They are placed here rather than in the main
body to maintain narrative focus.

\subsection{Geospatial Burden Distribution}

\begin{figure}[h]
  \centering
  \includegraphics[width=\linewidth]{india_fig2_spatial_maps}
  \caption{State-level spatial distribution of COVID-19 burden (GeoPandas
    static maps). Left: cases per lakh population (YlOrRd). Centre: CFR
    (\%, OrRd)---a clear North-South gradient visible. Right: vaccine
    doses administered (YlGn). The CFR gradient motivates the
    healthcare-infrastructure hypothesis discussed in
    Section~\ref{subsec:composite}.}
  \Description{Three choropleth maps of India, one per metric.}
  \label{app:fig:spatial}
\end{figure}

\subsection{State-Level Centrality Analysis}
\label{app:centrality}

\begin{figure}[h]
  \centering
  \includegraphics[width=\linewidth]{india_fig4_centrality}
  \caption{Six-metric centrality analysis of the India state correlation
    network (degree, betweenness, closeness, eigenvector, PageRank, Katz).
    Karnataka, Tamil Nadu, and Maharashtra rank highest across multiple
    metrics, identifying them as national-level epidemiological hubs.}
  \Description{Six-panel bar chart figure, one per centrality metric.}
  \label{app:fig:cent}
\end{figure}

\subsection{Spatial Adjacency Network}
\label{app:spatial_net}

\begin{figure}[h]
  \centering
  \includegraphics[width=\linewidth]{india_fig7_spatial_network}
  \caption{Inter-state spatial adjacency network (GeoPandas Queen
    contiguity). Nodes at state centroids; edges connect states sharing
    a land border. Node size $\propto$ confirmed cases; colour encodes
    CFR. Used as the simulation graph for the SIR analysis in
    Section~\ref{subsec:sir}.}
  \Description{Map-embedded network with nodes at state geographic
    positions and lines connecting adjacent states.}
  \label{app:fig:spatial_net}
\end{figure}

\subsection{Pune SEIR Model Fit}
\label{app:seir}

\begin{figure}[h]
  \centering
  \includegraphics[width=\linewidth]{pune_fig2_seir_fit}
  \caption{SEIR model fit for Pune District. Left: best-fit curve
    ($R_0 = 5.95$, $\beta = 0.5949$) overlaid on observed weekly cases.
    The low $R^2 = 0.009$ reflects a structural mismatch between a
    single-wave SEIR and the two-wave observed series; Wave~1 fit is
    accurate. Right: compartment dynamics over 420 days.}
  \Description{Two-panel SEIR figure. Left: bar chart with overlaid
    model curve. Right: S/E/I/R compartment trajectories.}
  \label{app:fig:seir}
\end{figure}

\subsection{Pune Mobility Network Visualisation}
\label{app:network}

\begin{figure}[h]
  \centering
  \includegraphics[width=\linewidth]{pune_fig4_network}
  \caption{Pune District mobility network. Nodes at geographic coordinates;
    size $\propto$ population; colour $\propto$ composite hub score.
    Edge width $\propto \log(F_{ij})$. The spoke-and-hub topology with
    Pune City and Haveli at the centre is immediately apparent.}
  \Description{Geographic network plot with 15 nodes and weighted edges.
    Pune City and Haveli are large red central nodes.}
  \label{app:fig:network}
\end{figure}

\subsection{Pune Metapopulation Scenario Comparison}

\begin{figure}[h]
  \centering
  \includegraphics[width=\linewidth]{pune_fig6_scenarios}
  \caption{Metapopulation SEIR scenario comparison (four pathogen
    profiles). Key result: for all scenarios with $R_0 \geq 1.4$,
    final attack rate exceeds 93\%---confirming that peak prevalence,
    not attack rate, is the actionable management target.}
  \Description{Four-panel scenario comparison figure.}
  \label{app:fig:scenarios}
\end{figure}

\subsection{Pune Network Robustness}

\begin{figure}[h]
  \centering
  \includegraphics[width=\linewidth]{pune_fig9_robustness}
  \caption{Network robustness under targeted vs.\ random hub removal.
    D50 = 8 nodes under targeted removal vs.\ $>$12 nodes for random
    removal, confirming scale-free-like vulnerability to hub attack.}
  \Description{Two-panel robustness figure with giant component fraction
    and average shortest path length curves.}
  \label{app:fig:robust}
\end{figure}

\section{Figure Saving Instructions}
\label{app:savefig}

Add the following before each \texttt{plt.show()} in the notebooks:

\begin{verbatim}
plt.savefig("india_fig1_epidemic_curve.png",
            dpi=180, bbox_inches="tight",
            facecolor="white")
\end{verbatim}

Full filename mapping is provided in the header comment of this
\texttt{.tex} file. Place all saved PNG files in the same directory
as the \texttt{.tex} before compiling.

\section{Software Environment}
\label{app:env}

Python~3.10; key packages: \texttt{networkx==3.1},
\texttt{geopandas==0.13}, \texttt{python-louvain==0.16},
\texttt{scipy==1.11}, \texttt{folium==0.14},
\texttt{matplotlib==3.7}, \texttt{seaborn==0.12},
\texttt{scikit-learn==1.3}, \texttt{libpysal==4.7},
\texttt{esda==2.5}, \texttt{pymannkendall==1.4}.

\end{document}
\endinput
